\documentclass[journal]{IEEEtran}

\usepackage{cite}
\usepackage{amsmath,amssymb,amsfonts}
\usepackage{graphicx}
\usepackage{booktabs}
\usepackage{tabularx}
\usepackage{xcolor}
\usepackage[hidelinks]{hyperref}

% Research-draft annotations (disabled once revision requirements are resolved).
\newif\ifshowrevisionnotes
\showrevisionnotesfalse

\newcommand{\revisiontodo}[1]{%
  \ifshowrevisionnotes
  \par\smallskip
  \noindent{\color{red}\textbf{Revision requirement:} #1}
  \par\smallskip
  \fi
}

\newcommand{\modelname}{KWT-GQA}

\title{Weight Sharing and Structured Compression for
Small-Footprint Keyword-Spotting Transformers}

\author{Ali Zare%
\thanks{Ali Zare is with Dotin Research and Studies Center,
Datis Arian Qeshm (Dotin), Tehran, Iran
(e-mail: alizare84@gmail.com).}}

\begin{document}
\maketitle

\begin{abstract}
Reducing the number of stored weights does not necessarily reduce
the computation or latency of a keyword-spotting model.
This paper studies a small-footprint Transformer that combines
cross-layer weight sharing with grouped-query attention.
A width-48 encoder repeatedly applies one shared block, while a
structured compression procedure removes intermediate convolutional
channels, feed-forward units, and attention groups without changing
the residual-stream width.
On the 12-class Google Speech Commands v1 task, the tied model
achieves $94.15\pm0.33\%$ accuracy with 36,516 unique parameters
over five training seeds.
Its untied counterpart has 219,204 parameters and achieves
$94.79\pm0.46\%$ over three seeds.
These results describe an approximately sixfold reduction in unique
parameters at unchanged architectural computation, rather than a
paired estimate of the accuracy effect of sharing.
Validation-selected structured pruning followed by dynamic
quantization of linear layers yields a 29,861-parameter model
with $93.83\pm0.41\%$ accuracy over five seeds.
Training the same reduced architecture from scratch and subsequently
quantizing it reaches $93.66\pm0.56\%$; the available results do not
establish an accuracy advantage for prune-then-fine-tune.
We distinguish unique parameters, arithmetic operations, tensor
payload, serialized artifact size, and measured runtime.
Analytical dense and pruned MAC totals are 31.273M and 24.071M,
respectively; tying alone does not reduce MACs.
The current evidence establishes a compact Transformer operating
point, but does not establish superiority over parameter-matched
untied Transformers or compact convolutional baselines.
\end{abstract}

\begin{IEEEkeywords}
Keyword spotting, Transformer, weight sharing,
grouped-query attention, structured pruning,
dynamic quantization.
\end{IEEEkeywords}

\section{Introduction}

Keyword spotting (KWS) identifies predefined words in short audio
segments and is commonly used as an entry point to speech interfaces.
Small-footprint KWS models must balance recognition accuracy with
storage, computation, and deployment constraints
\cite{zhang2017hello,warden2018speech}.
These constraints are related but not interchangeable.
A model may store few weights while repeatedly applying them,
and an integer-weight representation may reduce storage without
improving latency in a particular execution environment.

Compact convolutional models, including DS-CNN, TC-ResNet,
and BC-ResNet, provide strong reference points for this setting
\cite{zhang2017hello,choi2019temporal,kim2021bcresnet}.
Transformer-based KWS models demonstrate the applicability of
self-attention to short speech inputs
\cite{berg2021kwt}.
The question considered here is not whether a Transformer is
universally preferable to a convolutional network.
Instead, we examine how a Transformer can allocate a small
unique-parameter budget through cross-layer sharing and a reduced
number of key/value projections.

Cross-layer weight sharing offers a direct storage--computation
trade-off.
Reusing one block across multiple encoder iterations reduces the
number of independently stored weights, but retains the computation
of each iteration.
Grouped-query attention (GQA) reduces the number of key/value
projections while preserving multiple query heads
\cite{ainslie2023gqa}.
It also exposes attention groups as a structured removal unit.
These properties motivate studying sharing and GQA jointly,
without assuming that either improves accuracy under every budget.

We investigate these choices using \modelname{}, a KWT-style model
with a convolutional front-end, a width-48 encoder, and twelve
applications of a shared pre-normalized attention block.
After training, a structured compression procedure reduces
intermediate dimensions while keeping residual additions valid.
Dynamic quantization is then applied to linear layers as a separate
representation change.

The present study extends our earlier untied GQA-based KWS model
\cite{zareprev}.
The additional elements are cross-layer sharing, physical
reconstruction of structurally pruned tensors, and a multi-seed
evaluation on GSC v1-12.
The earlier v2 result is not used as a directly matched accuracy
baseline for the present v1 experiments.

Our current contributions are:
\begin{enumerate}
  \item An empirical characterization of a tied-GQA KWS model
  with 36,516 unique parameters, including a descriptive comparison
  with its untied counterpart.
  \item A structured compression procedure that preserves
  cross-layer sharing and residual compatibility, evaluated against
  training the same reduced architecture from scratch.
  \item A stage-separated evaluation framework that distinguishes
  architectural reductions from representation and runtime effects.
\end{enumerate}

The scope of these contributions is deliberately limited.
The available experiments do not establish that sharing is better
than conventional width or depth reduction at a matched parameter
budget, nor do they isolate the accuracy contribution of GQA from
tying: parameter-matched untied, shared-depth, and tied-MHA
controls are left to future work.
They also do not establish superiority over compact CNNs under a
common training and deployment protocol.

\section{Related Work}

\subsection{Small-Footprint Keyword Spotting}

Google Speech Commands provides a widely used benchmark for
small-footprint speech classification
\cite{warden2018speech}.
DS-CNN established practical convolutional baselines for
resource-constrained KWS
\cite{zhang2017hello}.
TC-ResNet and BC-ResNet subsequently explored efficient temporal
and broadcasted residual structures
\cite{choi2019temporal,kim2021bcresnet}.

These models are relevant reference architectures for the parameter
regime considered here.
However, published results may differ in dataset version,
class construction, preprocessing, augmentation, and training budget.
Consequently, literature values are not treated as matched
experimental evidence.

\subsection{Transformers, Grouped Queries, and Weight Sharing}

KWT applies Transformer machinery to keyword spotting
\cite{berg2021kwt}.
GQA reduces the number of key/value heads relative to query heads
\cite{ainslie2023gqa}.
For a fixed query-head count and sequence length, this reduction
does not remove the full per-query-head attention computation.

ALBERT-style parameter sharing reuses parameters across depth
\cite{lan2020albert}.
The relevant effect in the present setting is a reduction in unique
stored parameters rather than an automatic reduction in arithmetic
operations.
Our architecture combines this sharing mechanism with GQA in a
small-width KWS encoder.

\subsection{Pruning and Quantization}

Unstructured magnitude pruning removes individual nonzero weights
without necessarily changing dense tensor shapes
\cite{han2015learning}.
Structured pruning removes complete units or channels and can
therefore reduce dense architectural computation when the
corresponding tensors are rebuilt
\cite{li2017pruning}.

Integer quantization changes the representation and execution of
selected operations
\cite{jacob2018quantization}.
The dynamic quantization path used here is narrower than a
full-integer deployment graph: linear weights are quantized,
whereas several other model components remain floating point.

The novelty sought in this study is not a new general-purpose
pruning or quantization algorithm.
Rather, it is an empirical understanding of how established
mechanisms interact in a small tied-GQA KWS model.

\section{Research Questions and Evaluation Quantities}

\subsection{Research Questions}

The experimental design is organized around four questions:
\begin{enumerate}
  \item What storage--accuracy trade-off is obtained by sharing an
  encoder block rather than storing separate blocks?
  \item How do shared depth and the number of key/value groups
  affect accuracy and computational cost?
  \item Does prune-then-fine-tune provide an advantage over training
  the resulting small architecture directly?
  \item Which architectural and representation reductions
  translate into lower measured deployment costs?
\end{enumerate}

The existing results provide partial evidence for the first,
third, and fourth questions.
The second question and the parameter-matched interpretation of
the first require additional controls.

\subsection{Definitions}

We use the following distinct quantities.

\paragraph{Unique parameters.}
Let $\Theta_{\mathrm{unique}}$ contain each learned parameter tensor
once, regardless of how many times it is referenced or executed:
\begin{equation}
P_{\mathrm{unique}}
=
\sum_{\theta\in\Theta_{\mathrm{unique}}}
\operatorname{numel}(\theta).
\end{equation}
Quantized packed weights and learned biases remain part of this
count.
Non-learned buffers are reported separately.

\paragraph{Architectural arithmetic.}
MACs count multiply--accumulate operations for one model forward
pass on a specified feature tensor.
Where FLOPs are reported as $2\times\mathrm{MACs}$, this convention
is stated explicitly.
Softmax, normalization, nonlinearities, indexing, and quantization
overheads must either be counted separately or identified as
excluded.

\paragraph{Tensor payload.}
Payload is the byte count of unique stored tensor data, including
the stated quantization metadata and buffers.
The deduplication rule must identify the actual shared data,
rather than counting every module reference.

\paragraph{Serialized artifact size.}
Artifact size is the actual byte size of a saved, reloadable
representation.
It may exceed the tensor payload because of serialization metadata,
packing, and container overhead.

\paragraph{Runtime cost.}
Latency and peak memory depend on the hardware, runtime,
thread configuration, operator implementation, and measurement
boundary.
They cannot be inferred from parameter counts alone.

We use bytes as the primary storage unit.
When converted, $\mathrm{KB}=10^3$ bytes and
$\mathrm{KiB}=2^{10}$ bytes.

\section{Tied-GQA Architecture}

\subsection{Input and Convolutional Front-End}

The input is an MFCC representation with $T=98$ frames and
40 coefficients per frame, produced by the
\texttt{kws\_streaming} TensorFlow pipeline used in this
repository.
Audio is resampled to 16\,kHz.
Features use an 80-bin mel filterbank, a 30\,ms window,
a 10\,ms hop, and a 40-coefficient DCT, with the upper mel edge
at 7.6\,kHz.
Training applies the pipeline's default augmentations
(time shift, resampling, background mixing, and SpecAugment).
Features are then normalized with per-coefficient mean and
standard deviation estimated on the training set.
Sequences shorter than 98 frames are zero-padded; longer
sequences are randomly cropped in training and center-cropped
in evaluation.
TensorFlow feature extraction is not fully seed-controlled;
Torch/NumPy random state and cropping are.

Two one-dimensional convolutional layers map channels
$40\rightarrow48\rightarrow48$.
The front-end uses kernel size 3, stride 1, and padding 1,
with GELU and batch normalization, preserving the temporal length.

A learned absolute positional embedding of shape
$1\times98\times48$ is added before the encoder.
The input dropout probability is 0.05.

\subsection{Shared Encoder}

The residual width is $d=48$ and the feed-forward hidden width
is $d_{\mathrm{ff}}=96$.
A pre-normalized encoder block applies
\begin{align}
x &\leftarrow x+
\operatorname{Dropout}\!\left(
\operatorname{GQA}(\operatorname{LN}_1(x))
\right),\\
x &\leftarrow x+
\operatorname{Dropout}\!\left(
\operatorname{FFN}(\operatorname{LN}_2(x))
\right).
\end{align}
The FFN uses GELU and the residual dropout probability is 0.1.

In the tied model, the same block is applied $L=12$ times.
In the untied model, each iteration has independent block
parameters.
The forward-pass structure and tensor dimensions are otherwise
held fixed in the described comparison.

A relative-position bias table of shape
$(2T-1)\times H=195\times8$ is shared across encoder iterations.
The model also retains the absolute positional embedding.
We adopt both mechanisms following common practice in recent
acoustic Transformers; absolute-only, relative-only, and combined
ablations are not reported here and are left to future work.
A final normalization, temporal mean pooling, and a linear
classifier produce twelve logits.

\subsection{Grouped-Query Attention}

The model uses $H=8$ query heads, $G=4$ key/value groups,
and head dimension $d_h=6$.
The projections produce
\begin{equation}
Q\in\mathbb{R}^{T\times Hd_h},
\qquad
K,V\in\mathbb{R}^{T\times Gd_h}.
\end{equation}
Each key/value group serves the associated query heads.

After pruning, the number of query heads and key/value groups may
decrease while $d_h$ and residual width $d$ remain fixed.
Thus the query projection maps $d\rightarrow H'd_h$ and the output
projection maps $H'd_h\rightarrow d$.
This preserves the shape required by residual addition.

\subsection{Reported Parameter Inventory}

Table~\ref{tab:inventory} records the dense tied inventory
verified against checkpoints
(\texttt{dense\_A3\_seed*} / \texttt{untied\_A1\_seed*}).

\begin{table}[t]
\caption{Reported dense tied parameter inventory.}
\label{tab:inventory}
\centering
\begin{tabular}{lr}
\toprule
Component & Unique parameters\\
\midrule
Convolutional front-end & 12,960\\
Absolute positional embedding & 4,704\\
Relative-position bias & 1,560\\
Shared encoder block & 16,608\\
Final LayerNorm & 96\\
Classifier & 588\\
\midrule
Total tied & 36,516\\
Total untied & 219,204\\
\bottomrule
\end{tabular}
\end{table}

For a shared block containing $P_{\mathrm{block}}$ parameters
and all remaining parameters totaling $P_{\mathrm{other}}$,
\begin{align}
P_{\mathrm{tied}}&=P_{\mathrm{other}}+P_{\mathrm{block}},\\
P_{\mathrm{untied}}&=P_{\mathrm{other}}+LP_{\mathrm{block}}.
\end{align}
The ratio is not generally equal to $L$, because the front-end,
position embeddings, and classifier are not replicated.

\subsection{Analytical Attention Cost}

For full attention, the two matrix products $QK^\top$ and $AV$
require
\begin{equation}
M_{\mathrm{attention}}
=
2LT^2Hd_h
\label{eq:attention-macs}
\end{equation}
MACs, excluding projections and other operators.

Under the described dense dimensions, this is
11,063,808 MACs.
For the pruned configuration with $H'=6$ and unchanged
$d_h=6$, it is 8,297,856 MACs.

Whole-model counts are obtained by the same shape arithmetic
over the front-end, all twelve encoder applications
(projections, attention matmuls, and FFN), and the classifier,
excluding softmax, LayerNorm, GELU, and indexing.
This yields 31.273M MACs (62.546M FLOPs under the
$2\times\mathrm{MACs}$ convention) for the dense tied model and
24.071M MACs (48.142M FLOPs) for the selected pruned shape.
An earlier draft reported 10.358M / 8.139M MACs from the
\texttt{thop} profiler, which does not instrument
\texttt{torch.matmul} and therefore omitted the attention
matmuls above.
Table~\ref{tab:macs} gives the component breakdown.

\begin{table}[t]
\caption{Analytical MAC breakdown for one forward pass
($T{=}98$). Softmax, normalization, and elementwise nonlinearities
are excluded. FLOPs use $2\times\mathrm{MACs}$.}
\label{tab:macs}
\centering
\footnotesize
\begin{tabular}{lrr}
\toprule
Component & Dense (M) & Pruned (M)\\
\midrule
Conv front-end & 1.242 & 0.983\\
Attention projections ($\times12$) & 8.129 & 6.097\\
Attention matmuls Q$K^\top$+AV ($\times12$) & 11.064 & 8.298\\
FFN ($\times12$) & 10.838 & 8.693\\
Classifier & $<$0.001 & $<$0.001\\
\midrule
Total MACs & 31.273 & 24.071\\
Total FLOPs & 62.546 & 48.142\\
\bottomrule
\end{tabular}
\end{table}

\section{Tying-Preserving Structured Compression}

\subsection{Admissible Pruning Axes}

Residual width $d$ is kept fixed.
The following dimensions may be reduced:
\begin{enumerate}
  \item intermediate channels between the two front-end
  convolutions;
  \item FFN hidden units;
  \item complete GQA groups and their associated query heads.
\end{enumerate}

Removing a unit requires consistent slicing of all dependent
tensors.
For convolutional channels, this includes the first convolution's
outputs, the second convolution's inputs, and the relevant batch
normalization parameters and running statistics.
For FFN units, corresponding first-layer rows and second-layer
columns are removed.
For attention groups, query, key/value, output-projection, and
relative-bias indexing must remain consistent.

\subsection{Importance Ranking and Reconstruction}

Units are ranked by bidirectional L1 magnitude and then physically
rebuilt as smaller modules.
Only one copy of a shared block is reconstructed, and all encoder
iterations continue to reference that copy.

For an intermediate convolutional channel, the score is the sum of
outgoing Conv1 magnitudes (and bias, if present), incoming Conv2
magnitudes, and the associated batch-normalization scale when
present.
For an FFN hidden unit, the score is the sum of the corresponding
row of the first linear layer and column of the second, plus the
first-layer bias when present.
For a GQA group, the score aggregates L1 magnitudes over the
associated rows/columns of $Q$, $K$, $V$, and the output
projection.
Kept indices are the top-$k$ scores in descending order; at least
one unit is retained on each axis.

This bidirectional scoring is the implementation used for all
journal results.
It differs from the conference description, which used
incoming-only convolutional/FFN scores.

\subsection{Validation-Based Amount Selection}

Candidate pruning amounts are
$\alpha\in\{0.10,0.15,0.20,0.30\}$.
Each candidate is initialized from a trained dense tied checkpoint
and fine-tuned for 2,000 AdamW steps at learning rate $10^{-4}$.

The selection rule retains the most compressive candidate whose
mean validation accuracy drop across five seeds does not exceed
0.5 percentage points.
This is a constraint on the observed mean, not a per-seed
guarantee or a statistical upper confidence bound.

The selected $\alpha=0.20$ configuration has 38 intermediate
convolutional channels, 77 FFN units, six query heads, and three
key/value groups.
The reported parameter count is 29,861.

\subsection{Dynamic Linear Quantization}

After selecting the pruned architecture, PyTorch dynamic
quantization is applied to \texttt{nn.Linear} modules using
the \texttt{fbgemm} backend and a per-channel weight configuration
\cite{paszke2019pytorch}.

Linear weights are represented as INT8, whereas linear biases,
convolutions, normalizations, and position-related tensors remain
floating point under this path.
At linear computation time, activations are quantized dynamically
for the integer operation.
This is therefore a mixed execution graph rather than a
full-integer deployment.

Quantization does not reduce the number of architectural weight
elements.
It can change their representation, serialized size, and execution
cost.
Those effects are measured separately from shape reductions caused
by pruning.

\section{Experimental Protocol}

\subsection{Dataset and Training}

Experiments use Google Speech Commands v1 with ten target keywords,
unknown, and silence
(\texttt{yes}, \texttt{no}, \texttt{up}, \texttt{down},
\texttt{left}, \texttt{right}, \texttt{on}, \texttt{off},
\texttt{stop}, \texttt{go}, plus \texttt{unknown} and
\texttt{silence}).
Official hash-based splits yield 51,094 training, 6,798 validation,
and 6,835 test utterances.
Out-of-vocabulary words map to \texttt{unknown}; silence clips are
drawn from the corpus background-noise folder through the
\texttt{kws\_streaming} processor.

Training uses AdamW
\cite{loshchilov2019decoupled},
learning rate $10^{-3}$, weight decay $5\times10^{-4}$,
label smoothing 0.1, and batch size 512.
The schedule uses a ten-epoch linear warm-up followed by
cosine decay over a nominal 23,000-step budget.
Early stopping monitors validation loss with patience 15 and
minimum delta $5\times10^{-4}$; the best validation checkpoint is
retained for test evaluation.
On the hardware used for the multi-seed runs, a full dense train
is on the order of a few GPU-hours; the 2,000-step pruning
fine-tune is on the order of ten minutes.

\subsection{Variants and Seeds}

We distinguish:
\begin{itemize}
  \item U: untied GQA model;
  \item T: dense tied GQA model;
  \item T-Q: T with dynamic linear quantization;
  \item T-P: structurally pruned and fine-tuned T, before
  quantization;
  \item T-P-Q: T-P after dynamic linear quantization;
  \item S-Q: the same reduced architecture trained from scratch
  and subsequently quantized.
\end{itemize}

U currently uses seeds $\{0,1,2\}$ because seeds 3 and 4 were not
trained for the untied baseline; all U--T comparisons therefore use
only the three shared seeds and are descriptive.
T, T-Q, T-P-Q, and S-Q use seeds $\{0,1,2,3,4\}$.
The pruning sweep is reported over the five dense-T seeds.
For S-Q, weights and buffers of the reduced architecture are
reinitialized before training; no parent T-P state is retained.
T-P-Q is obtained by quantizing the fine-tuned float pruned
checkpoint for the same seed.

\subsection{Statistical Reporting}

Accuracy is reported as mean and standard deviation across seeds.
Paired comparisons use the same documented seed-level pairing
and explicitly define the direction of the difference.

A non-significant test is not interpreted as proof of equivalence.
A formal equivalence claim would require a justified margin and
an appropriate equivalence analysis.

Clip-level macro FAR and FRR are one-vs-rest quantities derived
from the twelve-class confusion matrix.
They are not streaming false accepts per hour and do not replace
threshold-based wake-word evaluation.

\subsection{Runtime Protocol}

Host latency is measured with \texttt{scripts/bench\_latency\_v2.py}
on an x86\_64 CPU under PyTorch 2.0 with the \texttt{fbgemm}
backend, \texttt{torch.set\_num\_threads(1)}, batch size one,
\texttt{model.eval()}, and synthetic MFCC-shaped inputs
($1\times98\times40$).
Each session uses 100 warm-up forwards and 1,000 timed forwards
with \texttt{torch.no\_grad()} and \texttt{time.perf\_counter}.
We report median, interquartile range, and 95th percentile over
timed forwards for T, T-Q, and T-P-Q.
These are host-CPU forward timings, not streaming end-to-end,
MCU, or energy measurements.

\section{Results}

\subsection{Sharing and the Dense Operating Point}

Table~\ref{tab:accuracy} summarizes the currently reported results.
The untied and tied models achieve
$94.74\pm0.35\%$ and $94.15\pm0.33\%$, respectively, both over five seeds.
The unique-parameter ratio is approximately six, while the
described computation graph retains twelve block executions.

Because no parameter-matched untied
control is available, this comparison is descriptive.
It does not establish that sharing is the best way to allocate
a fixed small parameter budget.

\begin{table}[t]
\caption{GSC v1-12 accuracy.
Values are mean $\pm$ sample standard deviation across seeds.
T-P is the float model after the selected pruning amount and
fine-tuning; T-P-Q is the same checkpoint after dynamic linear
quantization.}
\label{tab:accuracy}
\centering
\footnotesize
\begin{tabular}{lrrl}
\toprule
Variant & Seeds & Parameters & Accuracy (\%)\\
\midrule
U & 5 & 219,204 & $94.74\pm0.35$\\
T & 5 & 36,516 & $94.15\pm0.33$\\
T-Q & 5 & 36,516 & $94.18\pm0.34$\\
T-P & 5 & 29,861 & $93.93\pm0.36$\\
T-P-Q & 5 & 29,861 & $93.83\pm0.41$\\
S-Q & 5 & 29,861 & $93.66\pm0.56$\\
\bottomrule
\end{tabular}
\end{table}

\subsection{Structured Pruning Sweep}

Table~\ref{tab:sweep} reports the available sweep results.
The $\alpha=0.30$ candidate exceeds the mean validation-drop
ceiling, whereas $\alpha=0.20$ satisfies it and yields the
smallest eligible architecture.

\begin{table}[t]
\caption{Reported structured pruning sweep after 2,000
fine-tuning steps. Accuracy values are for float models.
Validation drop is in percentage points.}
\label{tab:sweep}
\centering
\footnotesize
\begin{tabular}{crll}
\toprule
$\alpha$ & Parameters & Val.\ drop & Test acc.\ (\%)\\
\midrule
0.10 & 34,211 & $-0.31\pm0.70$ & $94.73\pm0.26$\\
0.15 & 31,147 & $0.20\pm1.07$ & $94.11\pm0.36$\\
0.20 & 29,861 & $0.12\pm0.82$ & $93.93\pm0.36$\\
0.30 & 27,823 & $0.90\pm1.02$ & $93.60\pm0.55$\\
\bottomrule
\end{tabular}
\end{table}

The relatively high accuracy at $\alpha=0.10$ does not isolate
a benefit of removing weights.
The additional optimization steps are a confounding factor.
A matched control that continues dense T for 2,000 steps without
pruning (\texttt{scripts/run\_t\_ft\_control.py}) is currently running;
results will be incorporated once all five seeds are complete.
Until then, isolating continued training from structural adaptation
remains a limitation of the current evidence.
Amount selection used only the mean validation-drop rule above;
test-set sweep accuracies in Table~\ref{tab:sweep} are reported
for completeness and were not used to choose $\alpha$.

\subsection{Pruning and Quantization Effects}

Paired comparisons are recomputed from the five shared seeds
in \texttt{results/per\_seed\_metrics.csv}.
Dense linear quantization changes mean accuracy from
$94.15\%$ to $94.18\%$.
The paired difference $\Delta_s=A_{\mathrm{T-Q},s}-A_{\mathrm{T},s}$
has mean $+0.03$ percentage points
(std.\ $0.02$, $t=3.16$, df $=4$, two-sided $p=0.034$,
95\% CI $[0.01,0.05]$).
The effect is statistically detectable at $\alpha=0.05$ but
negligible in magnitude.

For the selected reduced architecture, the float and quantized
means are $93.93\%$ and $93.83\%$.
Most of the T-to-T-P-Q gap therefore arises before quantization.
For
\begin{equation}
\Delta_s=A_{\mathrm{T-P-Q},s}-A_{\mathrm{T},s},
\end{equation}
the mean is $-0.32$ percentage points
(std.\ $0.37$, $t=-1.92$, df $=4$, two-sided $p=0.13$,
95\% CI $[-0.78,+0.14]$).
These results neither demonstrate equivalence nor provide strong
evidence for a nonzero mean difference at the available sample size.

\subsection{Training the Reduced Architecture from Scratch}

S-Q achieves $93.66\pm0.56\%$, compared with
$93.83\pm0.41\%$ for T-P-Q.
For $\Delta_s=A_{\mathrm{S-Q},s}-A_{\mathrm{T-P-Q},s}$,
the mean is $-0.18$ percentage points
(std.\ $0.64$, $t=-0.61$, df $=4$, two-sided $p=0.57$,
95\% CI $[-0.97,+0.62]$).
No accuracy advantage for prune-then-fine-tune is established by
these results.

If a trained dense parent is already available, 2,000-step
fine-tuning can provide a low incremental-cost route to a smaller
model.
This does not imply lower total cost from a cold start:
parent training (nominal 23k-step budget with early stopping),
the $\alpha$ sweep with 2,000 fine-tuning steps per candidate,
and scratch training of the reduced architecture all consume
resources.

\subsection{Host Runtime}

Table~\ref{tab:latency} summarizes the audited host-CPU forward
timings.
Median latency is 12.10\,ms for T, 13.16\,ms for T-Q, and
12.38\,ms for T-P-Q.
Dynamic INT8 therefore does not accelerate this workload on the
measured x86\_64 host; the T-Q median is about 9\% slower than T,
while T-P-Q remains within a few percent of T.
Quantized models show lower 95th-percentile latency
(16.7--17.3\,ms versus 22.8\,ms for T), suggesting more predictable
tails rather than lower median cost.
Hardware with native INT8 acceleration may behave differently.

\begin{table}[t]
\caption{Single-thread CPU forward latency (batch 1, 100 warm-up,
1,000 timed runs). Times in milliseconds.}
\label{tab:latency}
\centering
\footnotesize
\begin{tabular}{lrrr}
\toprule
Variant & Median & IQR & P95\\
\midrule
T & 12.10 & $[10.54,15.45]$ & 22.75\\
T-Q & 13.16 & $[12.51,14.10]$ & 16.69\\
T-P-Q & 12.38 & $[11.85,13.39]$ & 17.25\\
\bottomrule
\end{tabular}
\end{table}

\subsection{Storage and Operation Counts}

Analytical MACs are reported in Table~\ref{tab:macs}:
31.273M for dense T and 24.071M for the selected pruned shape
(about 23\% fewer architectural MACs), while tying alone does not
reduce MACs.
Serialized artifact sizes on disk are approximately 182\,KB for T,
113\,KB for T-Q, and 937\,KB for U
(\texttt{results/storage\_breakdown.csv}).
The T versus U file-size gap reflects shared-parameter values and
container behavior; a na\"ive \texttt{state\_dict} still lists
per-iteration keys, so unique-parameter count remains the primary
storage claim.
Quantized payloads separate INT8 weights, scale/zero-point
metadata, and remaining FP32 tensors; dynamic quantization reduces
file size without changing the architectural MAC count.

\section{Discussion}

\subsection{Weight Sharing as a Storage--Computation Trade-off}

The current results support the feasibility of a KWS Transformer
with few unique parameters and repeated shared computation.
They do not imply that repeated computation is inexpensive.
Whether this trade-off is attractive depends on which resource
is constrained.

A parameter-matched shallow or narrow untied Transformer may
provide a different accuracy--latency balance.
Shared-depth controls are therefore necessary to determine whether
twelve iterations provide useful additional computation compared
with three or six iterations.

\subsection{The Role of GQA}

GQA reduces key/value projection dimensions and provides grouped
units for structured removal.
However, for fixed query-head count, head dimension, and sequence
length, the two attention matrix products retain the cost in
Eq.~\eqref{eq:attention-macs}.

The current experiments do not isolate GQA from tying.
A tied-MHA control is needed at fixed width, followed by a
separate comparison at a similar parameter budget.

\subsection{Pruning as Adaptation Versus Architecture Selection}

The scratch result suggests a useful distinction.
Pruning may serve as a way to adapt an existing trained model,
or as a way to discover a smaller architecture.
The present data do not establish that inherited weights improve
final accuracy over direct training of the selected architecture.

Accordingly, the contribution is not framed as a universal
accuracy benefit of pruning.
Its practical value must be assessed using adaptation cost,
final accuracy, and deployment cost together.

\subsection{Memory Beyond Stored Weights}

Small weight storage does not guarantee small working memory.
For example, an explicitly materialized FP32 attention tensor
with shape $8\times98\times98$ contains 307,328 bytes.
This is an illustrative tensor size, not a measured peak-memory
result.
Actual memory depends on operator fusion, allocation lifetimes,
and the execution backend.

This distinction motivates measuring runtime memory rather than
using parameter bytes as its proxy.

\section{Limitations}

The principal evidence currently uses GSC v1-12.
All five untied seeds are now complete and are included in the U row
of Table~\ref{tab:accuracy}.
A parameter-matched
untied depth-1 control (U-L1; 36{,}516 unique parameters), a tied-MHA
control ($G{=}H{=}8$), a matched T$+$2{,}000-step fine-tuning control
without pruning, and seed-0 position ablations remain queued under the
same training recipe
(\texttt{scripts/run\_revision\_controls.sh}).
Architecture and position controls therefore remain incomplete.

Compact CNNs (DS-CNN, TC-ResNet, BC-ResNet) have not been
retrained under the same recipe.
Literature CNN figures are therefore kept out of the main results
table and are not treated as matched experimental evidence.

The available runtime measurements are host-CPU forward timings,
not streaming end-to-end or energy measurements; peak memory was
not profiled.
ONNX exports of dense tied and untied checkpoints are provided for
offline MCU tooling (\texttt{results/onnx/}), but
STM32Cube.AI / ST Edge AI Analyze estimates (Flash, RAM, MAC/cycle
on Cortex-M4/M7) were not obtained in this environment because the
vendor tool is not installed here.
On-silicon measured latency and energy on a NUCLEO board remain
future work.

Dynamic quantization affects linear layers only.
The model is not a full-integer graph.
Operation counts are analytical MAC totals under the stated
exclusions; they are not profiler traces of every kernel.

Legacy noise results from a different pruning amount and an
unmatched checkpoint are excluded from the main evidence.
A valid robustness comparison requires matched parent checkpoints,
noise realizations, and final model configurations.

The small number of training seeds limits the precision of
between-method comparisons.
Non-rejection of a difference is not evidence of equivalence.

\section{Conclusion}

This study examines cross-layer weight sharing and structured
compression in a small-footprint KWS Transformer.
The reported tied model uses 36,516 unique parameters and reaches
$94.15\pm0.33\%$ accuracy on GSC v1-12 over five seeds.
Its untied counterpart stores approximately six times as many
parameters and reaches $94.74\pm0.35\%$ across five seeds,
a higher descriptive mean over the same number of seeds.

Structured pruning and dynamic linear quantization produce a
29,861-parameter model at $93.83\pm0.41\%$.
Training the same reduced architecture from scratch yields a
nearby mean accuracy, without evidence that prune-then-fine-tune
has an accuracy advantage.

The central implication is a storage--computation trade-off,
not a claim of universal Transformer efficiency.
Establishing when sharing and GQA are preferable requires
parameter-matched architecture controls.
Audited MAC totals, storage sizes, and host latency are reported
above; stronger deployment claims would need matched embedded
targets and energy measurements.

\section*{Reproducibility}

Training and evaluation scripts, model definitions, measurement
utilities, and machine-readable tables are intended for release at
\url{https://github.com/alizare84/LightTransformerKWS2-}.
At submission time, confirm that the public tree includes the
revision scripts and \texttt{results/*.csv} files cited here; large
checkpoints and the Speech Commands cache remain excluded
(\texttt{ARTIFACT\_RELEASE.md}).
A Zenodo DOI for a frozen release tag should be added once the
queued revision experiments are merged.

\section*{Acknowledgments}

Generative AI tools (large language models) were used to assist with
drafting and copy-editing of manuscript text and with scaffolding of
analysis scripts.
They were not used to invent experimental results, and they are not
authors of this work.
All numerical claims, tables, and scientific conclusions were
checked by the author against logs and checkpoints.
At submission time, any additional venue-specific AI-use declaration
(cover letter checkbox or dedicated statement) will follow the
target journal's current policy.

\section*{Compliance with Ethical Standards}

This study uses the publicly available Google Speech Commands
corpus and collects no new human-subject data.
Any additional declarations should follow the requirements of
the target journal and the author's institution.

\bibliographystyle{IEEEtran}
\bibliography{refs}

\end{document}